\documentclass[aspectratio=169]{beamer}

\usepackage[utf8]{inputenc}
\usepackage[T1]{fontenc}
\usepackage{amsmath}
\usepackage[
backend=biber,
style=bwl-FU,
sorting=ynt
]{biblatex}
\usepackage{graphicx}
%\usepackage[final]{changes}
\usepackage[draft]{changes}

\usepackage{tikz}
\usetikzlibrary{decorations.pathreplacing}

\addbibresource{references.bib} % Specify the bibliography file

\title{Financial Model: Traditional vs. Blockchain Storage}
\date{2024-10-16}
%\author{Malte Kerl}

\usetheme[language=english]{LTU}


\begin{document}
	
	\begin{frame}
		\titlepage
	\end{frame}
	
	\begin{frame}{Context}
		\begin{itemize}
        \item \textbf{Objective}: Develop a financial model to evaluate and compare traditional and blockchain-based data storage systems.
        \item \textbf{Motivation}:
        \begin{itemize}
            \item Increasing data storage needs demand efficient and secure solutions.
            \item Blockchain offers a decentralized alternative with unique costs and benefits.
        \end{itemize}
        \item \textbf{Goal}: Highlight trade-offs in costs, risks, and scalability.
    \end{itemize}
        
	\end{frame}
 
	\begin{frame}{Conceptual Framework}
    \begin{itemize}
        \item \textbf{Traditional Data Storage}
        \begin{itemize}
            \item Centralized infrastructure with single entity control.
            \item Key Costs: Operational, security, redundancy.
        \end{itemize}
        \item \textbf{Blockchain-Based Data Storage}
        \begin{itemize}
            \item Decentralized peer-to-peer network.
            \item Key Costs: Consensus, security, latency.
        \end{itemize}
        \item \textbf{Comparison Basis}: Financial sustainability, risk, and stability.
    \end{itemize}
    \end{frame}
 
    \begin{frame}{Utility Function Overview}
    \begin{itemize}
        \item \textbf{Purpose}: To analyze the adoption of each system.
        \item \textbf{Structure of Utility Functions}:
        \begin{itemize}
            \item Revenue/benefits minus system-specific costs.
            \item Incorporates network effects, risk, and stability considerations.
        \end{itemize}
        \item \textbf{Inspired by \cite{jalan_incentive-aware_2024} }: Incentive-aware agents, network interactions, cost-sharing.
    \end{itemize}
    \end{frame}

    \begin{frame}{Similarities between data storage and financial transfers}{Network Structure and Connectivity}
    
    Both financial and data storage networks involve strategic interactions among nodes. \vspace{3mm}
    
    \begin{itemize}
        \item \textbf{Financial Networks}: Banks or institutions form connections (loans or asset sharing) based on mutual benefits and risks.
        \item \textbf{Data Storage Networks (Blockchain)}: Nodes maintain data and achieve consensus with other nodes, validating each other’s data integrity.
        \item The network structure in both systems depends on the cost-benefit analysis of forming connections.
    \end{itemize}
    \end{frame}
    
    \begin{frame}{Similarities between data storage and financial transfers}{Network Stability and Resilience}
        
    Network stability in both financial and data storage networks relies on maintaining beneficial connections.
    \vspace{3mm}
    
    \begin{itemize}
        \item \textbf{Financial Networks}: Stability depends on sustaining links that provide risk-sharing or diversification benefits.
        \item \textbf{Blockchain Storage}: Stability is achieved when nodes are incentivized to stay active, preserving data integrity and availability.
        \item Stability in both networks emerges from balancing the cost of maintaining links with the overall benefit they provide.
    \end{itemize}
    \end{frame}
    
    % Slide 4: Cost Optimization and Incentives
    \begin{frame}{Similarities between data storage and financial transfers}{Cost Optimization and Incentives}
        
    Both models optimize utility by balancing the costs of participation with incentives or returns.
    \vspace{3mm}
    
    \begin{itemize}
         \item \textbf{Financial Networks}: Agents connect based on a cost-benefit analysis of maintaining financial relationships.
         \item \textbf{Blockchain Storage}: Nodes participate in data validation based on rewards (e.g., transaction fees) versus consensus costs.
        \item Each system seeks to maximize net utility by optimizing connections and incentives, aligning with \cite{jalan_incentive-aware_2024} incentive-aware approach.
    \end{itemize}
    \end{frame}
    
    % Slide 5: Network Impact on Risk and Security
    \begin{frame}{Similarities between data storage and financial transfers}{Network Impact on Risk and Security}

\added{Narrar Ejemplo 5 de Jalan et alii interpretado para el caso}
    
    The structure of both networks influences the spread and containment of risks.
    \vspace{3mm}
    
    \begin{itemize}
        \item \textbf{Financial Networks}: Risk of contagion arises when institutions are interlinked, potentially spreading failures.
        \item \textbf{Blockchain Storage}: Security risks (e.g., 51\% attacks) threaten data integrity across the network if nodes fail to maintain consensus.
        \item Both networks manage collective risks by structuring connections that either limit or control the spread of adverse events.
    \end{itemize}
    \end{frame}
    
    % Slide 6: Incentives for Participation and Stability
    \begin{frame}{Similarities between data storage and financial transfers}{Incentives for Participation and Stability}
        
    Both networks rely on effective incentives to maintain active participation and stability.
    \vspace{3mm}
        
    \begin{itemize}
        \item \textbf{Financial Networks}: Banks or institutions maintain connections if the benefits (liquidity, risk-sharing) exceed the costs.
        \item \textbf{Blockchain Storage}: Nodes are incentivized through rewards to participate in validation, maintaining network stability.
        \item Both systems use incentives to ensure that nodes continue to participate, securing the network's overall functionality and resilience.
    \end{itemize}
    \end{frame}
    
    \begin{frame}{Traditional Data Storage Utility Function}
        \begin{block}{Utility Function}
        \[
        U_T(\overbrace{w_T}^\text{completar}, P) = w_T^T (\mu_T - P e_T) - \gamma_T \cdot w_T^T \Sigma_T w_T - \sum_{j \in N(T)} c_{Tj} w_{Tj}
        \]
        \end{block}
        \begin{itemize}
            \item \( w_T^T (\mu_T - P e_T) \): Expected returns after operational costs.
            \item \( \gamma_T \cdot w_T^T \Sigma_T w_T \): Risk-adjusted returns.
            \item \( \sum_{j \in N(T)} c_{Tj} w_{Tj} \): Cost of connections for redundancy and security.
        \end{itemize}
    \end{frame}
    
    \begin{frame}{Blockchain-Based Data Storage Utility Function}
        \begin{block}{Utility Function}
        \[
        U_B(w_B, P) = w_B^T (\mu_B - P e_B) - \gamma_B \cdot w_B^T \Sigma_B w_B - \sum_{j \in N(B)} c_{Bj} w_{Bj}
        \]
        \end{block}
        \begin{itemize}
            \item \( w_B^T (\mu_B - P e_B) \): Expected rewards and benefits of decentralized storage.
            \item \( \gamma_B \cdot w_B^T \Sigma_B w_B \): Risk-adjusted returns.
            \item \( \sum_{j \in N(B)} c_{Bj} w_{Bj} \): Cost of maintaining connections with other nodes.
        \end{itemize}
    \end{frame}
    
    \begin{frame}{Comparative Utility Model}
        \begin{block}{Utility Difference}
        \[
        \Delta U = U_B(w_B, P) - U_T(w_T, P)
        \]
        \end{block}
        \begin{itemize}
            \item \( \Delta U > 0 \): Blockchain storage is more financially viable.
            \item \( \Delta U < 0 \): Traditional storage is preferable.
        \end{itemize}
        \begin{itemize}
            \item \textbf{Purpose}: Quantify relative advantage under different conditions.
        \end{itemize}
    \end{frame}

    \begin{frame}{Key Trade-Offs for Comparison}
        \begin{itemize}
            \item \textbf{Security and Redundancy}
            \begin{itemize}
                \item Traditional: Costly redundancy systems.
                \item Blockchain: Decentralization provides redundancy, with consensus costs.
            \end{itemize}
            \item \textbf{Operational and Consensus Costs}
            \begin{itemize}
                \item Traditional: Stable operational costs.
                \item Blockchain: Variable consensus costs with network size.
            \end{itemize}
            \item \textbf{Latency and Scalability}
            \begin{itemize}
                \item Traditional: Low latency due to centralization.
                \item Blockchain: Higher latency, affected by consensus.
            \end{itemize}
        \end{itemize}
    \end{frame}
    
    \begin{frame}{Conclusion}
        \begin{itemize}
            \item Presented a financial model comparing traditional and blockchain-based storage.
            
            \item \textbf{Future Directions}: Multi-period analysis, dynamic network adjustments, regulatory impacts.
        \end{itemize}
    \end{frame}
    
	\begin{frame}[allowframebreaks]
		\frametitle{References}
		\printbibliography
	\end{frame}
\end{document}